% ---------------------------------------------------------------------------
% ACL / EMNLP / NAACL Conference Paper Template
%
% Usage:
%   1. Download official acl.sty from: https://acl-org.github.io/ACLPUB/
%   2. Place acl.sty and acl_natbib.bst in this directory
%   3. Compile: pdflatex template.tex && bibtex template && pdflatex template.tex x2
%
% Conference: ACL (Association for Computational Linguistics) / EMNLP / NAACL
% Page limit: 8 pages (long) / 4 pages (short) + references
% ---------------------------------------------------------------------------

\documentclass[11pt,a4paper]{article}

% --- UNCOMMENT when acl.sty is available ---
% \usepackage[review]{acl}
% \usepackage{times}
% \usepackage{latexsym}
% \usepackage[T1]{fontenc}
% \usepackage[utf8]{inputenc}
% \usepackage{microtype}
% \usepackage{inconsolata}

\usepackage[utf8]{inputenc}
\usepackage[T1]{fontenc}
\usepackage{graphicx}
\usepackage{amsmath,amssymb}
\usepackage{booktabs}
\usepackage{multirow}
\usepackage[pagebackref,breaklinks,colorlinks,citecolor=blue]{hyperref}
\usepackage[margin=1in]{geometry}
\usepackage{times}

% --- Paper Information ---
% \aclfinalcopy

\title{Your Paper Title Here}

\author{
  Author One${}^1$,
  Author Two${}^1$,
  Author Three${}^2$ \\
  \\
  ${}^1$Institution One \\
  ${}^2$Institution Two \\
  \texttt{\{email1, email2\}@institution.edu}
}

\begin{document}

\maketitle

% ===========================================================================
% ABSTRACT
% ===========================================================================
\begin{abstract}
  Your abstract here. ACL format: typically single paragraph, $\sim$200 words.
  NLP papers often emphasize: task definition, method approach, key innovations,
  and empirical results.

  Follow 2-6-2 structure adapted for NLP:
  2 background/gap,
  6 method (First/Second/Finally),
  2 results/implication.
  Target: $\sim$200 words (ACL limit is typically 200 words).
\end{abstract}

% ===========================================================================
% 1. INTRODUCTION
% ===========================================================================
\section{Introduction}

% 5-paragraph structure, NLP-flavored:
% Para 1: Task importance + real NLP applications
% Para 2: Current progress -- pretrained LMs, fine-tuning, prompting approaches
% Para 3: Remaining gap -- specific limitations of existing NLP approaches
% Para 4: Proposed method -- core mechanism teaser
% Para 5: Contributions -- 3-4 bullets with evidence

Our main contributions are:
\begin{enumerate}
  \item Contribution 1.
  \item Contribution 2.
  \item Contribution 3.
\end{enumerate}

% ===========================================================================
% 2. RELATED WORK
% ===========================================================================
\section{Related Work}

% NLP papers often organize by: task, method paradigm, or resource
\subsection{Task-Specific Approaches}
\subsection{Pretrained Language Models for [Task]}
\subsection{[Method Paradigm]}

% ===========================================================================
% 3. PROPOSED METHOD
% ===========================================================================
\section{Proposed Method}
\label{sec:method}

\subsection{Problem Formulation}
\label{sec:method-problem}

Given [input definition], our goal is to [task definition].

\subsection{Model Architecture}
\label{sec:method-arch}

\subsection{Training Objective}
\label{sec:method-loss}

\begin{equation}
  \mathcal{L} = \mathcal{L}_{\text{task}} + \lambda \mathcal{L}_{\text{aux}}
  \label{eq:loss}
\end{equation}

\subsection{Inference}
\label{sec:method-inference}

% ===========================================================================
% 4. EXPERIMENTAL SETUP
% ===========================================================================
\section{Experimental Setup}
\label{sec:setup}

\subsection{Datasets}
\label{sec:setup-data}

\begin{table}[t]
  \centering
  \caption{Dataset statistics.}
  \label{tab:datasets}
  \begin{tabular}{lrrr}
    \toprule
    \textbf{Dataset} & \textbf{Train} & \textbf{Dev} & \textbf{Test} \\
    \midrule
    Dataset 1 & 0,000 & 0,000 & 0,000 \\
    Dataset 2 & 0,000 & 0,000 & 0,000 \\
    \bottomrule
  \end{tabular}
\end{table}

\subsection{Baselines}
\label{sec:setup-baselines}

\subsection{Implementation Details}
\label{sec:setup-impl}

\subsection{Evaluation Metrics}
\label{sec:setup-metrics}

% ===========================================================================
% 5. RESULTS AND ANALYSIS
% ===========================================================================
\section{Results and Analysis}
\label{sec:results}

\subsection{Main Results}
\label{sec:results-main}

\begin{table}[t]
  \centering
  \caption{Main results. Best in \textbf{bold}, second-best \underline{underlined}.}
  \label{tab:main}
  \begin{tabular}{lcccc}
    \toprule
    \textbf{Method} & \textbf{DS1} & \textbf{DS2} & \textbf{DS3} & \textbf{Avg} \\
    \midrule
    Baseline 1 & 00.0 & 00.0 & 00.0 & 00.0 \\
    Baseline 2 & 00.0 & 00.0 & 00.0 & 00.0 \\
    \textbf{Ours} & \textbf{00.0} & \textbf{00.0} & \textbf{00.0} & \textbf{00.0} \\
    \bottomrule
  \end{tabular}
\end{table}

\subsection{Ablation Studies}
\label{sec:results-ablation}

\subsection{Analysis}
\label{sec:results-analysis}

% Include error analysis, qualitative examples, attention visualization
% NLP papers benefit from qualitative analysis more than CV papers

% ===========================================================================
% 6. CONCLUSION
% ===========================================================================
\section{Conclusion}
\label{sec:conclusion}

% ===========================================================================
% LIMITATIONS (required at ACL/EMNLP/NAACL)
% ===========================================================================
\section*{Limitations}

Discuss limitations: computational requirements, language coverage,
domain specificity, ethical considerations, etc.

% ===========================================================================
% ETHICS STATEMENT (required at ACL)
% ===========================================================================
\section*{Ethics Statement}

% ===========================================================================
% ACKNOWLEDGMENTS
% ===========================================================================
\section*{Acknowledgments}

% ===========================================================================
% REFERENCES
% ===========================================================================
{
  \small
  \bibliographystyle{acl_natbib}
  \bibliography{references}
}

\end{document}
